\section{Generators}
\label{sec:generators}

The basic notion of a \emph{generator} --- a function $g: I \to \Zp$
encoding the body coefficients of an SCF --- was introduced in
Section~\ref{sec:scf} together with the three primitive types
\texttt{Generator}, \texttt{CachedGenerator}, and \texttt{FiniteGenerator}.
The present section focuses on the structural operations that transform
one generator into another: \texttt{advance}, \texttt{insert}, and
\texttt{prepend}.  It also describes \texttt{PeriodicGenerator}, which
extends the type hierarchy to eventually-periodic sequences.

\subsection{Structural operations}
\label{subsec:gen-ops}

All four generator types support three generator-to-generator operations.
Their specifications are given uniformly here; the per-type notes below
record any behavioural differences.

\begin{definition}[Advance]
  \label{def:advance}
  Let $g: \N \to \Zp$ be a generator and $n \geq 0$ an integer.  The
  \emph{$n$-advance} of $g$ is the generator
  \[
    \mathrm{adv}_n(g)(k) \;=\; g(n + k), \qquad k \geq 0.
  \]
  For $n = 0$, $\mathrm{adv}_0(g) = g$ (identity).
\end{definition}

\begin{remark}
  In terms of the SCF, advancing the generator by $n$ discards the first $n$
  body coefficients, which is precisely the operation performed by
  \texttt{SimpleContinuedFraction.tail(n)}: the $n$-th tail
  $[a_{n+1};\,a_{n+2},\ldots]$ has body generator
  $\mathrm{adv}_n(g)$.  The generator-level operation is therefore the
  foundational primitive underlying \texttt{tail}.
\end{remark}

\begin{definition}[Insert]
  \label{def:insert}
  Let $g: \N \to \Zp$ and let $f = (b_0, \ldots, b_{s-1})$ be a
  \texttt{FiniteGenerator} of size $s$.  For $\mathtt{at} \geq 0$, the
  \emph{insertion of $f$ at position $\mathtt{at}$} is the generator
  \[
    \mathrm{ins}(g,f,\mathtt{at})(k) =
    \begin{cases}
      g(k) & k < \mathtt{at}, \\
      b_{k - \mathtt{at}} & \mathtt{at} \leq k < \mathtt{at} + s, \\
      g(k - s) & k \geq \mathtt{at} + s.
    \end{cases}
  \]
\end{definition}

\begin{definition}[Prepend]
  \label{def:prepend}
  The \emph{prepend} of $f$ onto $g$ is $\mathrm{ins}(g, f, 0)$: the
  sequence $b_0, \ldots, b_{s-1}, g(0), g(1), \ldots$
\end{definition}

\begin{remark}[Prepend and the inverse operation]
  Proposition~\ref{prop:inverse} (case $a_0 > 0$) shows that inverting
  a positive SCF amounts to prepending the singleton $[a_0]$ to the body
  generator.  Conversely, case $a_0 = 0$ removes the first body coefficient
  to form the new integer part.  Both are special cases of advance/prepend on
  the generator.
\end{remark}

\begin{catenadef}[\texttt{advance}, \texttt{insert}, \texttt{prepend}]
\begin{lstlisting}[style=python]
from catena.generators import Generator, FiniteGenerator

g = Generator(lambda n: n + 1)   # 1, 2, 3, 4, ...
f = FiniteGenerator([10, 20])

# advance: shift the index by 3 -> 4, 5, 6, ...
g3 = g.advance(3)
[g3(k) for k in range(4)]        # [4, 5, 6, 7]

# insert f at position 2: 1, 2, 10, 20, 3, 4, ...
h = g.insert(f, at=2)
[h(k) for k in range(6)]         # [1, 2, 10, 20, 3, 4]

# prepend: 10, 20, 1, 2, 3, ...
p = g.prepend(f)
[p(k) for k in range(5)]         # [10, 20, 1, 2, 3]
\end{lstlisting}
\end{catenadef}

\paragraph{Type-specific behaviour.}

\begin{description}[leftmargin=2em, style=nextline]
  \item[\texttt{Generator}]
    All three operations return a plain \texttt{Generator} wrapping a
    new closure; no memoisation is carried over.

  \item[\texttt{CachedGenerator}]
    \texttt{advance} and \texttt{insert} accept an optional
    \texttt{copy\_cache: bool} flag (default \texttt{False}).
    When \texttt{True}, cache entries are transferred to the new generator
    with indices adjusted to reflect the transformation: for an advance by
    $n$, entry $(k, v)$ with $k \geq n$ is copied as $(k - n, v)$; for an
    insert at position $p$ with block size $s$, entries before $p$ are
    copied unchanged and entries from $p$ onwards are shifted to $k + s$.
    When \texttt{False} the new generator starts with an empty cache,
    which is appropriate when the closure captured by the new generator
    differs from the original (e.g., it calls back into the original generator
    rather than the underlying function directly).

  \item[\texttt{FiniteGenerator}]
    \texttt{advance}$(n)$ returns the suffix $g[n:]$ as a new
    \texttt{FiniteGenerator}; advancing past the end raises
    \texttt{IndexError}.  \texttt{insert} concatenates the underlying
    arrays, upgrading to the smallest unsigned typecode that covers all
    values in the combined sequence.  There is no \texttt{prepend} override;
    it delegates to \texttt{insert(fg, at=0)}.

  \item[\texttt{PeriodicGenerator}]
    See Section~\ref{subsec:periodic-gen}.
\end{description}

\subsection{PeriodicGenerator}
\label{subsec:periodic-gen}

A \texttt{PeriodicGenerator} encodes an eventually-periodic sequence: a
finite pre-period $(d_0, \ldots, d_{m-1})$ followed by an infinitely
repeating period $(c_0, \ldots, c_{r-1})$:
\[
  g(k) =
  \begin{cases}
    d_k & 0 \leq k < m, \\
    c_{(k - m) \bmod r} & k \geq m.
  \end{cases}
\]
It is the generator type used internally by
\texttt{PeriodicSimpleContinuedFraction}.

\paragraph{Advance.}
Let $n \geq 0$.
\begin{itemize}
  \item If $n < m$: the new pre-period is $(d_n, \ldots, d_{m-1})$; the
        period is unchanged.
  \item If $n \geq m$: the pre-period is empty and the period is rotated by
        $(n - m) \bmod r$ positions:
        \[
          (c_r, c_{r+1}, \ldots, c_{r-1}) \mapsto
          (c_{r'}, c_{r'+1}, \ldots, c_{r'-1})
          \quad\text{where } r' = (n - m) \bmod r.
        \]
\end{itemize}
In both cases the result is a \texttt{PeriodicGenerator} with the same or
shorter pre-period and a period that is a rotation of the original.

\paragraph{Insert.}
Inserting a \texttt{FiniteGenerator} $f$ of size $s$ at position
$\mathtt{at}$:
\begin{itemize}
  \item If $\mathtt{at} \leq m$: $f$ is spliced into the pre-period;
        the period is unchanged.
  \item If $\mathtt{at} > m$: one copy of the period is unrolled into
        the pre-period up to the insertion point, $f$ is spliced in, and
        the original period is kept.  Concretely, the new pre-period is
        \[
          (d_0, \ldots, d_{m-1},\,
           c_0, \ldots, c_{s'-1},\,
           b_0, \ldots, b_{s-1},\,
           c_{s'}, \ldots, c_{r-1})
          \quad\text{where } s' = (\mathtt{at} - m) \bmod r,
        \]
        and the period remains $(c_0, \ldots, c_{r-1})$.
\end{itemize}

\begin{remark}
  After an insert with $\mathtt{at} > m$, the period is
  \emph{not} lengthened; only one instance of the period is unrolled into the
  pre-period to accommodate the inserted block.  The resulting sequence is
  still eventually periodic with the same period, just a longer pre-period.
  This is the minimal representation: a shorter pre-period would require
  the inserted block to appear infinitely often, which is not the case.
\end{remark}

\begin{catenadef}[\texttt{PeriodicGenerator.advance} and \texttt{insert}]
\begin{lstlisting}[style=python]
from catena.generators import FiniteGenerator, PeriodicGenerator

# period (1, 2, 3), pre-period (9, 8)
pg = PeriodicGenerator(period=[1, 2, 3], pre_period=[9, 8])
# sequence: 9, 8, | 1, 2, 3, 1, 2, 3, ...

# advance by 1: pre-period shrinks to (8,)
pg.advance(1).pre_period   # FiniteGenerator([8])
pg.advance(1).period       # FiniteGenerator([1, 2, 3])

# advance by 3 (past the pre-period): period rotates by (3-2)%3 = 1
pg.advance(3).pre_period   # FiniteGenerator([])
pg.advance(3).period       # FiniteGenerator([2, 3, 1])

# insert [7] at position 4 (within the first copy of the period):
# pre-period becomes (9, 8, 1, 2, 7, 3), period stays (1, 2, 3)
ins = pg.insert(FiniteGenerator([7]), at=4)
ins.pre_period             # FiniteGenerator([9, 8, 1, 2, 7, 3])
ins.period                 # FiniteGenerator([1, 2, 3])
[ins(k) for k in range(10)]
# [9, 8, 1, 2, 7, 3, 1, 2, 3, 1]
\end{lstlisting}
\end{catenadef}

\paragraph{Quadratic coefficients.}
\texttt{PeriodicGenerator} also exposes
\texttt{quadratic\_coefficients()}, which returns the primitive polynomial
$(A, B, C)$ satisfied by the value of the generator (with integer part
zero).  The derivation is described in
Section~\ref{subsubsec:pscf-to-surd}; it is not repeated here.
